\chapter{Models of the primes}\label{ch:oc-cramer}

This chapter compares Cram\'er's random model of the primes with the theorems it explains and with Maier's theorem, which shows where the model fails.

\section{Cram\'er's model}

Cram\'er's model declares the events ``$n$ is prime'', for $n\ge3$, to be independent, with probability $1/\log n$. The expected number of primes up to $x$ is then asymptotic to $\operatorname{li}(x)$.

\section{Successes of the model}

The model predicts the typical size of the gap after a prime near $x$, of order $(\log x)^2$ at the extreme, and it predicts Poisson behaviour for primes in suitably random intervals.

\section{Maier's theorem}

Maier (1985) proved that the asymptotic $\pi(x+y)-\pi(x)\sim y/\log x$ fails for a sequence of $x$ when the length $y$ is a fixed power of $\log x$. The primes are not independent at that scale.

\section{Granville's refinement}

Granville's refinement thins the model by the prime factors of the primorial near $\log x$. The aim is to see which predictions survive the refinement and which were artefacts of independence.
